\documentclass[10pt,a4paper]{amsart}
\usepackage[margin=19mm]{geometry}
\usepackage{amsmath,amssymb,mathtools}
\usepackage[scr=rsfs,cal=euler]{mathalfa}
\usepackage{xfrac}
\usepackage[unicode,hidelinks,bookmarks=false]{hyperref}
\newcommand{\ie}{i.e.\ }
\newcommand{\cf}{cf.\ }
\newcommand{\resp}{respectively\ }
\newcommand{\Subsection}{\subsection}
\providecommand{\nobreakdash}{\nobreak\hbox{-}}
\setlength{\emergencystretch}{2em}
\multlinegap0pt
\usepackage{amssymb}
\usepackage{mathtools}
\usepackage{xcolor}
\usepackage[shortlabels]{enumitem}
\newtheorem{theorem}{Theorem}
\newcommand{\HH}{\mathcal{H}}
\newcommand{\JJ}{\mathcal{J}}
\newcommand{\T}{\mathbb{T}}
\newcommand{\abs}[1]{\lvert#1\rvert}
\newcommand{\ddmu}{\,d\mu}
\title{\textbf{A nonlocal transmission problem on a hybrid continuous-discrete domain}}
\author{Hafida Abbas, Abdelhalim Azzouz}
\date{Source: arXiv:2603.14178v1 (CC BY 4.0)}
\begin{document}
\maketitle

\noindent\emph{Source and licence.} \textbf{A nonlocal transmission problem on a hybrid continuous-discrete domain}. Version arXiv:2603.14178v1 (CC BY 4.0), distributed by arXiv under the Creative Commons Attribution 4.0 International licence, http://creativecommons.org/licenses/by/4.0/, as recorded in the arXiv licence metadata for this exact version and shown on the abstract page https://arxiv.org/abs/2603.14178v1. Full licence notice: \texttt{licenses/arxiv-2603.14178-CC-BY-4.0.txt}.\par\smallskip

\noindent\textit{Editorial scope.} Counted unit: the article's theorem thm:weak-EL (source0.tex lines 502--514). The article's complete original proof of it is delivered with it (source0.tex lines 516--584, one \texttt{proof} environment of 69 lines). No mathematics is written, repaired or re-derived: the statement and the proof are the article's own lines, in the article's own notation, and no retained line is substituted or re-flowed.  Retained uncounted as the source-local context the proof needs: lines 402--405; lines 471--473.  Nothing was omitted from any retained span, so the package carries no omission record.  No retained line contains a citation, so the wrapper carries no bibliography and no bibliography entry.  No retained line contains a figure environment, an image-inclusion command or drawing code, so no image is delivered and the package declares no asset dependency.  Editorial formatting only, in addition: an AMS \texttt{amsart} preamble that restates the article's own declarations and macros used by the retained lines, namely the environments \texttt{theorem} on separate counters, together with the standard packages those lines need.  The retained environments print in this document as Theorem 1 in order of appearance, whereas the article prints the same items with the numbering of its own full document.  The complete native source of the article as distributed in the arXiv e-print for this exact version is bundled unchanged as \texttt{sources/196382/source0.tex}.
\begin{equation}\label{eq:weak-solution}
\mathfrak a(u,\varphi)+\lambda\int_{\T}u\varphi\ddmu=\int_{\T}f\varphi\ddmu
\qquad\text{for every }\varphi\in\HH_\alpha(\T).
\end{equation}

\begin{theorem}[Variational characterization of weak solutions]\label{thm:weak-var-equivalence}
A function $u\in\HH_\alpha(\T)$ is the unique minimizer of $\JJ$ if and only if it is the unique weak solution of \eqref{eq:weak-solution}.
\end{theorem}

\begin{theorem}[Weak Euler--Lagrange identity]\label{thm:weak-EL}
Let $u=(v,a_2,\dots,a_{N+1})\in\HH_\alpha(\T)$ be the minimizer of $\JJ$. Then, for every $\varphi=(\psi,b_2,\dots,b_{N+1})\in\HH_\alpha(\T)$,
\begin{align}\label{eq:split-weak-form}
&\int_0^1\int_0^1
\frac{(v(x)-v(y))(\psi(x)-\psi(y))}{\abs{x-y}^{1+2\alpha}}\,dx\,dy \notag\\
&\quad +2\sum_{k=2}^{N+1}\int_0^1
\frac{(a_k-v(x))(b_k-\psi(x))}{\abs{k-x}^{1+2\alpha}}\,dx \\
&\quad +\sum_{\substack{i,j=2\\ i\neq j}}^{N+1}
\frac{(a_i-a_j)(b_i-b_j)}{\abs{i-j}^{1+2\alpha}}
+\lambda\int_{\T}u\varphi\ddmu
=\int_{\T}f\varphi\ddmu.\notag
\end{align}
\end{theorem}

\begin{proof}
Let $u=(v,a_2,\dots,a_{N+1})\in\HH_\alpha(\T)$ be the minimizer of $\JJ$. By Theorem~\ref{thm:weak-var-equivalence}, $u$ is the unique weak solution of \eqref{eq:weak-solution}. Hence, for every test function
\[
\varphi=(\psi,b_2,\dots,b_{N+1})\in\HH_\alpha(\T),
\]
one has
\begin{equation}\label{eq:proof-thm71-weak}
\mathfrak a(u,\varphi)+\lambda\int_{\T}u\varphi\,\ddmu=\int_{\T}f\varphi\,\ddmu.
\end{equation}
It therefore remains to compute the bilinear form $\mathfrak a(u,\varphi)$ in terms of the continuous and discrete components.

By definition,
\[
\mathfrak a(u,\varphi)=\iint_{\T\times\T}
\frac{(u(x)-u(y))(\varphi(x)-\varphi(y))}{\abs{x-y}^{1+2\alpha}}\,d\mu(x)\,d\mu(y).
\]
Using the decomposition of the measure
\[
d\mu(x)=\mathbf 1_{[0,1]}(x)\,dx+\sum_{k=2}^{N+1}\delta_k(x),
\]
the product measure $\mu\otimes\mu$ splits into four contributions:
continuous--continuous, continuous--discrete, discrete--continuous, and discrete--discrete. We examine them separately.

For the continuous--continuous part, since $u|_{[0,1]}=v$ and $\varphi|_{[0,1]}=\psi$, we obtain
\[
I_{cc}:=\int_0^1\int_0^1
\frac{(v(x)-v(y))(\psi(x)-\psi(y))}{\abs{x-y}^{1+2\alpha}}\,dx\,dy.
\]
For the continuous--discrete part, using $u(k)=a_k$ and $\varphi(k)=b_k$, we get
\[
I_{cd}:=\sum_{k=2}^{N+1}\int_0^1
\frac{(v(x)-a_k)(\psi(x)-b_k)}{\abs{x-k}^{1+2\alpha}}\,dx.
\]
For the discrete--continuous part, by the symmetry of the kernel and the identity $\abs{k-x}=\abs{x-k}$, we have
\[
I_{dc}:=\sum_{k=2}^{N+1}\int_0^1
\frac{(a_k-v(x))(b_k-\psi(x))}{\abs{k-x}^{1+2\alpha}}\,dx.
\]
Since
\[
(v(x)-a_k)(\psi(x)-b_k)=(a_k-v(x))(b_k-\psi(x)),
\]
it follows that $I_{cd}=I_{dc}$. Therefore the mixed contribution equals
\[
I_{cd}+I_{dc}=2\sum_{k=2}^{N+1}\int_0^1
\frac{(a_k-v(x))(b_k-\psi(x))}{\abs{k-x}^{1+2\alpha}}\,dx.
\]
Finally, for the discrete--discrete part we obtain
\[
I_{dd}:=\sum_{i=2}^{N+1}\sum_{j=2}^{N+1}
\frac{(a_i-a_j)(b_i-b_j)}{\abs{i-j}^{1+2\alpha}}.
\]
The terms with $i=j$ vanish because $(a_i-a_i)(b_i-b_i)=0$, hence
\[
I_{dd}=\sum_{\substack{i,j=2\\ i\neq j}}^{N+1}
\frac{(a_i-a_j)(b_i-b_j)}{\abs{i-j}^{1+2\alpha}}.
\]
Combining the four pieces yields
\begin{align*}
\mathfrak a(u,\varphi)
&=\int_0^1\int_0^1
\frac{(v(x)-v(y))(\psi(x)-\psi(y))}{\abs{x-y}^{1+2\alpha}}\,dx\,dy\\
&\quad+2\sum_{k=2}^{N+1}\int_0^1
\frac{(a_k-v(x))(b_k-\psi(x))}{\abs{k-x}^{1+2\alpha}}\,dx\\
&\quad+\sum_{\substack{i,j=2\\ i\neq j}}^{N+1}
\frac{(a_i-a_j)(b_i-b_j)}{\abs{i-j}^{1+2\alpha}}.
\end{align*}
Substituting this identity into \eqref{eq:proof-thm71-weak} gives exactly \eqref{eq:split-weak-form}. This completes the proof.
\end{proof}

\end{document}
